\section{Real-data Experimental Protocol}
We evaluate Graph-CoLD on two verified real datasets: CICIDS-2017 postfilter11 and CESNET-TLS-Year22 postfilter25. CICIDS-2017 uses the full postfilter11 protocol after removing classes below the minimum-count threshold and downsampling the dominant class. CESNET-TLS-Year22 uses a deterministic audit-window subset followed by postfilter25 stratified splitting; we do not describe it as a full-archive evaluation.

Each result row records source verification, dataset hash, sample policy, split id, noise seed, model seed, and active views. CICIDS-2017 uses host, IP, and temporal views. CESNET-TLS-Year22 uses IP and temporal views; process and threat-intelligence views are not claimed for this dataset.

Noise models include clean labels, symmetric label noise, asymmetric label noise, and graph-consistency noise. Noise is injected only into training labels. Graph-consistency rows use the beta values present in the D5.5 result matrix, and the beta=0 rows are retained to verify consistency with symmetric corruption.

The formal baselines are CoLD, ablation_hard, Noisy-Supervised, Confident-Learning, and Co-Teaching-lite. FINE, MCRe, MORSE, Flash, Argus, Decoupling, and full Co-Teaching are excluded from the formal comparison because the repository does not contain independently implemented and smoke-passed real-data rows for them. We report Macro-F1, FPR, FNR, ERR, Tail-ERR, compression ratio, runtime, and memory. Statistical tests are paired by dataset, noise type, noise rate, graph beta, and seed, avoiding unpaired pooled tests.
